% New border-rank findings. Include in the checked-results part only after root review.
\section{Arithmetic replacement of the seventh-minimum certificates}

\subsection*{BR-001: an independent finite checker for all 25 prolongation cases}
Let the six pairs of $\{0,1,2,3\}$ be in lexicographic order, and number the
rectangles $q=(I,J)$ by $6\operatorname{index}(I)+\operatorname{index}(J)$.
For a rectangle subset $S$, set $A=k^S$ and let $B=C$ have the used variables
$x_{rc}$ as basis.  Its residual inclusion tensor is
\[
 T_S^\epsilon=\sum_{q=((r,s),(c,d))\in S}a_q\otimes
 (x_{rc}\otimes x_{sd}+x_{sd}\otimes x_{rc}
 +\epsilon x_{rd}\otimes x_{sc}+\epsilon x_{sc}\otimes x_{rd}),
 \qquad \epsilon\in\{-1,1\}.
\]
Here $\epsilon=-1$ is the determinant and $\epsilon=1$ the permanent.
In characteristic zero, the following 25 cases have border rank at least 14.  The listed determinant
cases are also taken with their transposes.
\[
\begin{array}{c|c|l}
\text{inherited locator}&\epsilon&S\\\hline
\text{P-09 A}&-1&0,1,3,6,7,9,12\\
\text{P-09 B}&-1&0,1,3,6,7,12,13\\
\text{P-12 C}&-1&0,1,2,3,6,7,8\\
\text{P-12 D}&-1&0,1,2,3,4,6,12\\
\text{P-14 E}&-1&0,1,2,3,6,7,9\\
\text{P-14 F}&-1&0,1,2,3,6,7,18\\
\text{P-15 det}&-1&0,1,3,6,7,9,18\\\hline
\text{P-13 1}&1&0,1,2,3,4,6,12\\
\text{P-13 2}&1&0,1,2,3,6,7,8\\
\text{P-13 3}&1&0,1,2,6,7,8,18\\
\text{P-13 4}&1&0,1,2,6,7,18,20\\
\text{P-13 5}&1&0,1,3,6,7,9,14\\
\text{P-16 1}&1&0,1,2,3,6,7,9\\
\text{P-16 2}&1&0,1,2,3,6,7,18\\
\text{P-16 3}&1&0,1,3,6,7,15,21\\
\text{P-15 per A}&1&0,1,3,6,7,9,18\\
\text{P-15 per B}&1&0,1,3,6,7,18,21\\
\text{P-19}&1&0,1,2,6,7,18,19
\end{array}
\]
The characteristic-zero values were claims of BR\_PENCIL.  The new part here
is their independent reconstruction by standard-library programs and their
arithmetic extension, used in BR-002 below.

\paragraph{Proof and complete algorithm.}
We state the characteristic-independent form of the necessary condition,
since the inherited findings stated it only in characteristic zero.
Write $U=T(C^*)\subset A\otimes B$.  If $\underline R(T)\le r$, there is
an $r$-plane $E\supset U$ such that
\[
 \dim\ker\Psi_A(E)\ge r,\qquad \dim\ker\Psi_B(E)\ge r,
\]
where $\Psi_A:E\otimes A\to B\otimes\Lambda^2A$ and
$\Psi_B:E\otimes B\to A\otimes\Lambda^2B$ are wedge maps.
For spans of $r$ independent pure $AB$ tensors, the corresponding symmetric
prolongation tensors give $r$ independent kernel vectors; contraction by a
linear form nonzero on all factors proves independence.  The source bundles
have constant rank on the Grassmannian, so these kernel-dimension conditions
are closed.  The condition $U\subset E$ is closed and the Grassmannian is
proper.  Projection therefore passes this necessary condition to all border
limits.  The open condition that the $r$ pure $AB$ tensors are independent
is dense in the $r$-summand parameter space.  This proves the assertion over
any algebraically closed field.  This is the all-field form of BR\_PENCIL
P-09's first reduction and Conner--Huang--Landsberg, Proposition 2.1.

The split torus gives $x_{rc}$ weight $R_r+C_c$ and $a_q$ weight
$-\sum_{r\in I}R_r-\sum_{c\in J}C_c$.  A separating cocharacter limit in the
proper incidence permits $E$ to be torus invariant.  The image $U$ is the
direct sum of its $\dim B$ distinct weight lines, each with coefficients
$\pm1$.  In each such weight, omit the last nonzero coordinate of its image
vector.  The remaining coordinate vectors, together with $U$, are a basis
of $A\otimes B$ over $\mathbb Z$.  This specifies all quotient coordinates
without an arithmetic exceptional prime.

For each orientation and target weight, the programs form the wedge-map
columns explicitly and eliminate the image of $U$.  If $V_w$ is a quotient
weight block and $x\in V_w$, the projected column indexed by a factor basis
vector $f_i$ is a linear map $C_{w,i}x$ of target weight $w+\operatorname{wt}(f_i)$.
These target weights are distinct for fixed $w$.  Thus adjoining a line
$kx$ increases rank by precisely the number of its nonzero projected columns.
All the $C_{w,i}$ are recomputed, rather than read from the inherited data.

The finite enumeration is as follows.
\begin{enumerate}
\item Enumerate the support of the quotient $E/U$, respecting the capacity
of every weight block.  For quotient dimension one or two, the supports
are single blocks of sufficient capacity or pairs of distinct blocks.
For quotient dimension four, enumerate every support of at most four blocks
whose capacities sum to at least four.
\item Saturate each selected block by its whole quotient weight space.
If the kernel of either wedge map on this larger space has dimension
less than $r$, the support is impossible: the kernel for $E$ injects into it.
\item For a quotient line, enumerate every possible simultaneous zero
pattern of the $C_{w,i}$.  Start with no equations, repeatedly adjoin the
row space of one $C_{w,i}$, and keep intersections with a nonzero kernel.
Canonical exact row reduction removes repetitions.  An intersection $K$
represents the stratum obtained by removing the kernels of the remaining
maps that do not vanish identically on $K$.  It is nonempty over the infinite
algebraic closure: a finite union of proper linear subspaces cannot cover
$K$.  Conversely, every vector's exact zero pattern occurs in this enumeration.
\item For lines in distinct weights, the number of distinct active target
weights is a joint rank lower bound.  Another valid bound sums the individual
increments and subtracts the total possible overlap of target weights.
For a fixed line, quotient by its actual image before testing the other line.
\item For the support-nine cases the only surviving multiplicity patterns
are $(1,1,1,1)$ and $(2,1,1)$.  In the first pattern, quotient by the images
of every fixed line.  One or two parameter lines remain.  Repeat the exact
arrangement enumeration in this quotient and retain a refined stratum only
if every formerly required nonzero map remains nonzero somewhere on it.
For a plane in a three-dimensional block, parameterize it as $\ker n$.
A column map of rank one restricts to rank zero exactly when $n$ spans its
row space; a map of rank two restricts to rank one exactly when $n$ annihilates
its one-dimensional kernel; ranks zero and three restrict to ranks zero and
two.  These are linear conditions on $n$, so the same complete arrangement
algorithm applies.  Sum, over target weights, the largest contributed rank.
\end{enumerate}
These operations exhaust all subspaces, including the parameter families.
No generic sample stands in for a projective weight space.

For the first 20 cases, \texttt{check\_twodim\_extensions.py} gives:
\[
\begin{array}{c|r|r|r|r}
&\text{supports}&\text{saturated}&\text{line tests}&\text{conditional}\\\hline
 P09A&1095&1042&53&0\\
 P09B&50&47&3&0\\
 P12C&52&50&2&0\\
 P12D&60&58&2&0\\
 P14E&1189&1185&4&0\\
 P14F&1143&1076&66&1\\
 P13.1&60&58&2&0\\
 P13.2&52&49&3&0\\
 P13.3&50&45&5&0\\
 P13.4&50&46&4&0\\
 P13.5&56&56&0&0\\
 P16.1&1189&1176&13&0\\
 P16.2&1143&1075&67&1\\
 P16.3&1190&1173&17&0
\end{array}
\]
The six determinant representatives account for twelve cases by transposition.
The two support-nine permanent types and determinant type are handled by
\texttt{check\_support9.py}; the determinant transpose is the fourth case.
Their total supports are respectively $46119,51985,46119$, with no survivors.
The exact counts, including all parameter strata, are in the output JSON files.

For the remaining P-19 case, the first two steps leave eight supports.
Exact joint refinement gives eight fixed extensions, each of dimension 13.
The higher condition is
\[
 \dim\mathcal P_{22}(E)\ge r,
 \quad
 \mathcal P_{22}(E)=
 (A\otimes E\otimes B)\cap
 ((A^{\otimes2})^{\mathfrak S_2}\otimes(B^{\otimes2})^{\mathfrak S_2}).
\]
It is again a fixed-ambient kernel condition and therefore closed.  Its
necessity follows from the $r$ independent tensors $u_i^{\otimes2}\otimes
v_i^{\otimes2}$ for a span of pure tensors, with the same proper-incidence
argument.  Using invariant tensors, including diagonal entries, is essential
in characteristic two; no division by two occurs.
The program \texttt{check\_higher\_prolongation.py} first computes the kernel
of the $385\times143$ $B$-wedge map: rank 130 and kernel dimension 13.
It embeds this kernel in $A\otimes B^{\otimes2}$, verifies symmetry by exact
coefficient equality, and forms the next $A$-wedge map.  That map has domain
dimension 91 and rank 79.  Thus $\dim\mathcal P_{22}(E)=12$ in all eight cases,
a contradiction.  These computations reconstruct the matrices from $S$;
they import no old certificates.  This completes the independent proof.

\subsection*{BR-002: forty-three in every characteristic}
For the integral divided-power tensors fixed in BR\_PENCIL FL-01 and every
algebraically closed field $k$,
\[
 \mu_7^{\det}=\mu_7^{\mathrm{per}}=14,
 \qquad
 \underline R_k(S_{\det,4})\ge43,
 \qquad
 \underline R_k(S_{\mathrm{per},4})\ge43.
\]
The assertion includes characteristics two and three.

\paragraph{Proof.}
Use the field-independent coordinate degeneration BR\_PENCIL P-01 and
its determinant compression P-07.  These are integral root-subgroup
specializations, so P-07 remains valid in characteristic two.  There are
25 determinant Borel seven-plane ideals.  For the permanent, the 1517
low-support coordinate orbits under row and column permutations and
transposition cover exactly 1543608 subsets; support at least 14 gives
a unit flattening minor.  These finite cover assertions, carried from
BR\_PENCIL R-08/R-10, are also independently rerun by
\texttt{check\_cover.py}, which exhausts all $\binom{36}{7}=8347680$ subsets.

For every remaining non-prolongation case, form the entire Koszul map,
rather than the inherited chosen minor.  If its exterior factor has dimension
$t$ and degree $p$, a pure tensor has map rank $\binom{t-1}{p}$ in every
characteristic: choose its nonzero factor as a basis vector and inspect the
basis wedges.  Rank is subadditive and its upper-bound locus is closed.
Consequently map rank $\rho$ gives border rank at least
$\lceil\rho/\binom{t-1}{p}\rceil$.
The independent checker \texttt{check\_koszul.py} forms all columns directly
from $T_S^\epsilon$, separates exact torus-character blocks, and computes
all block ranks over $\mathbb Q$.  All 11 determinant support/Koszul cases
and 1506 permanent Koszul cases give at least 14.

Here is the arithmetic justification, beyond a computation at one large prime.
Every rational elimination records the numerator and denominator of every
nonunit pivot.  Direct nonzero scalar tests in the arrangement-compatibility
step also record their numerators and denominators.  The primes that occur
divide six.  Every row operation and every inverse used by the proof is
therefore defined over $\mathbb Z[1/6]$; every exact zero identity holds there.
All rank computations, linear intersections, and their quotient maps commute
with specialization to characteristics other than two and three.  A nonzero
branch test also remains nonzero there.  Intersections that become identical
can only duplicate already enumerated strata, not create omitted ones.
This proves the full calculation uniformly away from two and three.

The exceptional computations are small and separate.
\begin{itemize}
\item In characteristic three, all eleven determinant support/Koszul cases
again give at least 14.  The two-dimensional-extension tests have no prime
three in their pivots.  The support-nine determinant arrangement is rerun
over $\mathbb F_3$ and has no survivors.  The permanent Koszul calculation
has only prime two in its pivots; its support-nine cases are also rerun over
$\mathbb F_3$ with no survivors.
\item In characteristic two, determinant and permanent tensors coincide.
The determinant compression reduces to the same 25 Borel ideals.
All eleven support/Koszul cases still give at least 14 (the two support-factor
ranks become 6457, still greater than $13\cdot495=6435$).
Five of the six ordinary determinant prolongation representatives pass the
same exhaustive computation over $\mathbb F_2$.  The P-09 A representative
has eight survivors of the first-stage test; their joint arrangements give
eight fixed extensions.  The higher checker gives $\mathcal P_{22}$ dimension
12 for each, exactly as above.  Its transpose follows by swapping rows and
columns.  The support-nine determinant tensor equals the permanent A tensor,
whose full enumeration over $\mathbb F_2$ has no survivors.
\end{itemize}
Thus every required seventh residual tensor has border rank at least 14
in every characteristic.  The integral anchored-shape curve BR\_PENCIL P-06
supplies the matching upper bound.  Finally the flattening image of the
fourfold tensor has dimension 36, with unit coefficient generators, and
BR\_PENCIL R-01 gives $36-7+14=43$.

\paragraph{What is and is not claimed about unit minors.}
This proof supplies arithmetic replacements for the rank certificates;
it does not assert that all 9038 originally selected determinants are units.
In fact a literal replacement retaining every full block rank is impossible,
as BR-003 shows.  Rank losses in small characteristics are harmless when the
remaining total rank still exceeds the required pure-tensor threshold.

\subsection*{BR-003: a genuine nonunit block in the original certificate architecture}
For the determinant subset $S=\{0,1,2,3,6,7,12\}$, take the $B$-Koszul map
of exterior degree four.  In the weight
$(2,0,0,1;1,0,0,2)$, the complete nonzero block, in the ordered bases specified
by \texttt{even\_minor\_block.json}, is
\[
 M=\begin{pmatrix}
 1&0&1&0\\1&0&0&1\\0&1&0&0\\0&0&1&1
 \end{pmatrix}.
\]
It has determinant 2 and rank 3 in characteristic two.  In particular
there is no unit four-minor in that block, nor any four-minor nonzero in
characteristic two.  The original characteristic-zero full block rank cannot
be preserved by merely changing a selected minor in this block.
The complete support-factor rank nevertheless drops only from 6501 to 6457,
which still proves the same border-rank lower bound 14.
The proof is the displayed determinant computation and the independent full
block reconstruction; its first three independent rows certify rank three
modulo two.

\subsection*{Algebraic-geometric dictionary for BR-001--BR-003}
All factors are finite free modules over $\mathbb Z$.  Wedge maps are morphisms
between their exterior-power bundles, and their rank conditions define
closed determinantal loci.  Candidate extension spaces form closed incidences
in relative Grassmannians.  Projectivity supplies proper projection and
cocharacter specialization.  The arithmetic result proves emptiness in each
algebraically closed fiber: first over $\operatorname{Spec}\mathbb Z[1/6]$
by invertible elimination, then separately over the fibers at two and three.
It does not infer emptiness of special fibers from emptiness of the generic
fiber.  Invariant symmetric tensor squares in the higher test are free
integral modules, so the proof also specifies the characteristic-two model.

The inherited exact foundational locators are EGA II, Definition 5.4.1,
Proposition 5.4.2, and Theorem 5.5.3(i) (properness and projective morphisms);
FGA, Bourbaki expos\'e 221, Theorems 3.1--3.2 (Quot/Hilbert parameter spaces);
SGA 3, Expos\'e I, Sections 4.5--4.6 including Definition 4.6.1 (module and
linear group functors); and SGA 6, Expos\'e V, Definition 2.1, Examples 2.2(b),
and Lemma 2.2.1 (exterior operations on $K_0$).  These are framework locators
carried from BR\_PENCIL's checked dictionary, not fresh facsimile inspections
and not attributions of the modern numerical bounds to Grothendieck.
No unpublished archive cote/scan for these findings has been established.
No purity, weight, $\ell$-adic, or motivic statement is used.

\subsection*{BR-004: the determinant ten-plane gate reduces from 35 to 23}
\textbf{Status: proved by the following argument and exact computations;
root full proof/code review and independent rerun completed.  Characteristic zero only.}

Let the six two-subsets of $\{0,1,2,3\}$ be ordered lexicographically,
and index a rectangle by $6i+j$, with row-pair index $i$ and column-pair
index $j$.  The partial order is componentwise order on the four entries
of the two pairs.  Among the 59 lower order ideals of cardinality ten,
the residual determinant tensor has border rank at least 18 for 36
ideals.  Consequently the characteristic-zero determinant gate
$\mu_{10}\ge18$ reduces to the 23 ideals represented, up to transpose,
by the following twelve rows.  This does not prove $\mu_{10}\ge18$.
\[
\begin{array}{l|r|r|l}
 S&m&\text{rigid full spaces}&\text{first parameter branches: depth/count}\\\hline
0,1,2,3,4,5,6,7,8,9&12&1&4/1\\
0,1,2,3,4,5,6,7,8,12&14&1&\text{none}\\
0,1,2,3,4,5,6,7,8,18&12&1&4/1\\
0,1,2,3,4,5,6,7,9,12&13&0&4/1\\
0,1,2,3,4,5,6,7,9,18&11&0&4/3,\ 5/2\\
0,1,2,3,4,6,7,8,9,10&12&0&4/2\\
0,1,2,3,4,6,7,8,9,18&12&0&4/1\\
0,1,2,3,4,6,7,8,18,19&12&0&4/1,\ 5/1\\
0,1,2,3,6,7,8,9,18,19&12&0&3/1,\ 4/8\\
0,1,2,3,6,7,8,18,19,20&12&0&4/1,\ 5/3\\
0,1,2,3,6,7,9,12,18,19&13&0&3/1,\ 4/6\\
0,1,2,3,6,7,9,18,19,21&11&0&2/2,\ 3/7,\ 4/22,\ 5/16
\end{array}
\]
The penultimate row is self-transpose; the others have distinct transposes.
Each of the 80 parameter branches in this table is a projective line at
the first nonrigid flag step.  It is a base for the remaining invariant-flag
incidence, not a claim that a full admissible seventeen-plane is already
parametrized by that projective line.  Branches can overlap, and no
irreducible-component count is asserted.  The full rational basis for every
weight space, every preceding flag, and each two-dimensional kernel is
recorded in \texttt{certificates/det10\_gate.json}.

\paragraph{Proof: the stronger necessary condition.}
Write $T\in A\otimes B\otimes C$, with $\dim A=10$,
$B=C$ the $m$-dimensional used-variable space, and $U=T(C^*)$.
Let $Y_r$ be the closure in $\operatorname{Gr}(r,A\otimes B)$ of
the spaces spanned by $r$ independent rank-one matrices.  If
$\underline R(T)\le r$, the proper incidence of pairs $(E,T)$ with
$E\in Y_r$ and $T\in E\otimes C$ shows that some $E\in Y_r$
contains $U$.  The incidence fiber is projective and is preserved by
the upper triangular row and column groups.  A Borel fixed point therefore
gives such an invariant $E$.

For every $E\in Y_r$, the tautological tensor
$\tau_E\in E^*\otimes A\otimes B$, representing the inclusion
$E\hookrightarrow A\otimes B$, satisfies $\underline R(\tau_E)\le r$.
Indeed, on a Grassmannian chart choose a regular frame for the tautological
bundle.  At a span of $r$ rank-one matrices its inclusion tensor is a sum
of $r$ pure tensors; rank equations are closed, so they persist in the
limit in this local frame.  Changes of frame do not change border rank.
This is also the full-extension principle in Conner--Huang--Landsberg,
arXiv:2009.11391v1, proof of Proposition 2.4.

Every subspace $D\subset E$ has a tautological tensor obtained by
restricting the $E^*$ factor of $\tau_E$, hence
$\underline R(\tau_D)\le r$.  For $X=A$ or $B$, the Koszul map
\[
 K_{X,p}(\tau_D):D\otimes\Lambda^pX
       \longrightarrow Y\otimes\Lambda^{p+1}X,
 \qquad \{X,Y\}=\{A,B\},
\]
has rank at most $r\binom{\dim X-1}{p}$, because a pure tensor
has exactly that Koszul rank.  Thus a violation for a partial extension
$D=U+F$ excludes every completion to $E$.  The familiar first wedge
test is $p=1$; BR-004 also uses $p=2,3$.

\paragraph{Exhaustive invariant-flag construction.}
For each of the 59 ideals form the rational $H$-module
$V=(A\otimes B)/U$, decomposed by torus characters.  Every invariant
$(17-m)$-plane $E/U$ has an invariant complete flag, since $H$ is
connected solvable in characteristic zero.  At a partial flag space $F$,
all possible next invariant lines of weight $w$ lie in
\[
 \mathbb P\ker\left((V/F)_w\longrightarrow
       \bigoplus_{\alpha}(V/F)_{w+\alpha}\right),
\]
where the six maps are the adjacent row and column root operators.
The operators are computed from the exterior action on the minor labels
and the ordinary variable action: $-R_\alpha^{\mathsf t}\otimes1+
1\otimes N_\alpha$.  The checker also verifies invariance of $U$.
There is no enumeration of mere coordinate subsets within a repeated
weight space: its entire common kernel is computed over $\mathbb Q$.
A zero kernel supplies no line; a one-dimensional kernel supplies its
unique line over every characteristic-zero extension; a larger kernel
is exported as a genuine projective family and left open.

Every generated partial space is subjected to the preceding Koszul
inequalities.  The first-wedge ranks are computed over $\mathbb Q$.
For $p=2,3$ they are computed over the proven prime $1000003$ after
inverting each explicitly checked rational denominator.  A nonzero
minor modulo this prime is a nonzero rational minor.  Thus these
computations are lower-rank certificates over characteristic zero,
not numerical estimates and not claims about every characteristic.

Induction on flag length proves completeness: any hypothetical invariant
$E$ either passes through one of the exported parameter branches or
passes through a rigid successor at every step.  In the latter case its
partial tensors cannot fail a necessary rank inequality.  All such paths
are therefore accounted for.  The higher-extension enumeration excludes
32 ideals (17 transpose representatives) and leaves 27.  It independently
recomputes all eighteen P-17 cases and also excludes twelve ideals that
were absent from the former 24-case gate.

\paragraph{The three projected-Koszul witnesses and the union count.}
Three additional classical projected-Koszul witnesses, each with its
transpose, are independently reconstructed using integral $5\times10$
matrices $P:A\to k^5$.  For exterior degree two their ranks are
\[
\begin{array}{l|r|r}
S&\rank K_{k^5,2}((P\otimes1\otimes1)T)&\text{pure cost}\\\hline
0,1,2,3,4,5,6,7,18,19&104&6\\
0,1,2,3,4,5,6,12,18,24&104&6\\
0,1,2,3,4,6,7,9,18,19&106&6
\end{array}
\]
All exceed $17\cdot6=102$.  The matrices, not inherited rank assertions,
are the input in \texttt{det10\_projection\_inputs.json}; the checker
recomputes their ranks.  The middle row and its transpose already belong
to the 32 higher-extension exclusions.  The union therefore has
$32+6-2=36$ distinct ideals, leaving exactly 23.  The full partition of
all 59 ideals is asserted by the checker, which generates the poset ideals
afresh.  The total twelve new exclusions relative to the previous 24-case
gate are listed explicitly in the companion manifest.

\paragraph{A particularly small surviving gate.}
For $S=\{0,1,2,3,4,5,6,7,8,12\}$, the first-wedge tests alone leave
five invariant seventeen-planes.  Four have $B$-Koszul degree-two ranks
$1335,1335,1335,1336$, exceeding $17\binom{13}{2}=1326$.
Only
\[
 E_*=U+\langle a_9\otimes x_{02},\ a_9\otimes x_{03},\
                    a_9\otimes x_{10}\rangle
\]
remains, where $a_9$ is the dual label of rectangle 12 in the ordered list
$S$.  Its degree-two $B$-Koszul rank is 1306 and is not an exclusion.
The same statement holds for the $AC$ extension by symmetry.  The supplied
complete-flag probe of this fixed $E_*$ leaves a flag passing the tested
necessary inequalities; it supplies neither an upper bound nor a lower
bound of eighteen.  Constructing a rank-seventeen degeneration of this
extension would, by restriction, give an upper bound seventeen for this
residual tensor; excluding its viability would eliminate this gate.

\paragraph{Checker contract and exact residual incidence.}
Run \texttt{python scripts/check\_det10\_gate.py} from the border work
directory.  It uses only the Python standard library and the new local
linear-algebra modules, imports no inherited project code, and reconstructs
all signed tensors, root operators, quotients, flag kernels, and map ranks.
The certificate includes every excluded partial flag and its rejecting map.

For the coupled search, let
$V=(C\otimes E_{AB})\cap(B\otimes E_{AC})$ and let $R\subset B\otimes C$
be the span of all $A$-contractions of $V$.  A necessary condition is the
existence of an invariant $G=E_{BC}\cap R$ with
\[
 T(A^*)\subset G\subset R,\quad \dim G\le17,\quad
 \underline R(\tau_G)\le17,\quad
 \dim(V\cap(A\otimes G))\ge17.
\]
The original exploratory search incorrectly restricted to $\dim G=17$.
Enlarging a smaller $G$ inside $R$ need not preserve its rank inequalities,
so that restriction is not used in any exclusion count.  The corrected
\texttt{probe\_bc\_contractions.py} tests all dimensions from 10 through
17 before extending, and explicitly leaves all parameter branches open.
For $E_{AB}=E_{AC}=E_*$ it has no rigid survivor but 50 first projective
branches; this remains an open coupled incidence.

\paragraph{Geometric dictionary.}
The proof uses proper incidence in a relative Grassmannian, its flag
refinement, torus-character subbundles, and degeneracy loci of exterior
bundle maps.  The foundational EGA/FGA/SGA locators above apply unchanged.
It uses no purity theorem and no motivic specialization assertion.

\subsection*{BR-005: exact pencils and projected extension maps leave one transpose pair}
\textbf{Status: exact checker passes; root full proof/code review and independent rerun completed.
Characteristic zero only.}
Of the 59 determinant ten-plane Borel ideals, 57 have residual border
rank at least eighteen.  The only remaining ideals are
\[
 S_\diamond=\{0,1,2,3,6,7,9,18,19,21\},\qquad
 S_\diamond^{\mathsf t}=\{0,1,3,6,7,9,12,18,19,21\}.
\]
The gate $\mu_{10}\ge18$ is therefore reduced to this single transpose
pair.  No upper bound seventeen for either tensor is asserted.

\paragraph{Proof of the exact pencil refinement.}
Use the complete flag cover of BR-004.  A first parameter branch has
the form $D(t)=D_0+\langle X+tY\rangle$ for finite $t$, together
with $D(\infty)=D_0+\langle Y\rangle$, where $X,Y$ have the same
torus character.  For any of the relevant Koszul maps, first eliminate
the columns supplied by $D_0$ over $\mathbb Q$.  The remaining map
is a direct sum by torus weight of pencils $M_w(t)=A_w+tB_w$.
In each block choose a square minor $d_w(t)$ that is nonzero as a
polynomial, of order $r_w$.  If the fixed rank plus $\sum_w r_w$
exceeds the pure-tensor threshold, a surviving parameter must be a root
of some $d_w$, or the point at infinity.  This is an exact finite
exception set; selecting a nonzero minor at an integer value is only
a way of finding its row and column indices.  The determinant polynomial
itself is then computed and factored over $\mathbb Q$.

At a root of an irreducible factor $f(t)$ the rank is computed in the
field $\mathbb Q[t]/(f)$, which gives the same rank at every conjugate
root.  All factors encountered have degree at most two.  The complete
set of distinct nonconstant factors is
\[
 t,\ t+1,\ t+\tfrac12,\ t+2,\ t-1,\ t-\tfrac13,
 \ t-2,\ t-\tfrac23,\ t^2-t-1.
\]
Thus no real-only argument, random parameter sample, or numerical root
approximation is being used.  Infinity is checked separately in every
case.  Multiple Koszul tests are intersected by evaluating the remaining
finite points in the next test.

This process excludes 66 of BR-004's 80 first projective branches.
Six branches reduce to explicit finite sets of rational points; eight
remain open.  The latter all belong to $S_\diamond$.
For example, the sole branch for
$S=\{0,1,2,3,4,5,6,7,9,12\}$ has $B$-wedge rank at least 206
at every finite parameter, above the allowed $17\cdot12=204$.
At infinity its $B$ rank is 204, but its $A$-wedge rank is
157, above $17\cdot9=153$.

\paragraph{The rigid survivors and finite specializations.}
All three rigid seventeen-planes left in BR-004, including $E_*$,
are excluded by an integral projection of the $B$ factor onto $k^7$,
followed by the degree-three Koszul map of their tautological tensors.
Each projected map has rank 346 modulo $1000003$, whereas a border-rank
seventeen tensor has rank at most
$17\binom{6}{3}=340$.  The exact projection matrices are retained in
\texttt{det10\_refined\_gate.json}; their ranks are recomputed by
\texttt{extension\_projection.py}.  This use of a projection differs
from applying a Koszul map only to the original residual tensor: it
uses the stronger necessary condition on the limiting extension.

The six finite branches are then continued by the same exhaustive
invariant-line algorithm, now at their exact rational parameter values.
Every further projective kernel is again subjected to the exact pencil
procedure; none is silently replaced by sample points.  The two support-12
representatives
$\{0,1,2,3,6,7,8,9,18,19\}$ and
$\{0,1,2,3,6,7,8,18,19,20\}$ are excluded at their first specialized
step by the projected $B$-Koszul test.  The other four finite branches,
all on $S_\diamond$, are excluded after at most two further flag steps.
The saved trace includes every excluded flag and its rejecting rank.

Combining the preceding cover with BR-004 proves that exactly 57 of
the 59 ideals have been excluded.  The eight currently open branches
of $S_\diamond$ are those with indices
$0,1,2,3,5,9,12,13$ in its explicit BR-004 parameter manifest,
at depths $2,2,3,3,3,4,4,4$, respectively.  They are bases of
continuation incidences, not a list of viable border decompositions.

\paragraph{Verification and dictionary.}
Regenerate BR-004 first, then run
\texttt{python scripts/check\_det10\_refinement.py}.
The checker uses only the Python standard library and SymPy.  It
reconstructs the 80 pencil maps and exact exceptional fibers, the
three projected-rank witnesses, and all six finite continuations.
Its final assertions are 66 excluded/6 finite/8 open first branches,
three ranks equal to 346, six excluded finite continuations, and a
57/2 partition of the 59 ideals.  The recorded run passed in 32.43
seconds.  The argument is the same proper Grassmannian/flag incidence
as BR-004, with determinantal loci pulled back to projective lines.
The exceptional-point analysis is algebraic over $\mathbb Q$; it
does not specialize from characteristic zero to infer a positive-
characteristic theorem.  No purity or motivic assertion is used.

\subsection*{BR-006: the determinant ten-plane lower gate closes, giving forty-four}
\textbf{Status: proved.  The restarted consolidated exact checker returned
PASS in 710.2006 seconds: six family nodes, four quadratic-field restarts,
empty pending queue, and no unhandled branch.  The root reviewed the whole
saved proof and every new or adapted source module, with no fault found,
and accepted the completed certificate.  An earlier run was interrupted
by an execution-environment outage and is not counted as a PASS.}

Over every algebraically closed field of characteristic zero,
\[
 \mu_{10}(U^{\mathrm{det}}_{4,2})\ge18,
 \qquad
 \underline R_{\mathrm{CP}}(\operatorname{Det}_4)\ge44.
\]
The latter is the polynomial coefficient tensor in $(k^{16})^{\otimes4}$.
Its symmetric border rank is consequently also at least 44.
This finding gives a lower bound for $\mu_{10}$, not an exact value:
no matching upper bound eighteen is established here.  The currently
proved upper bound is nineteen by BR-007 below.  No permanent ten-plane conclusion follows
from the determinant calculation.

\paragraph{Reduction to three projective families.}
BR-005 leaves only $S_\diamond$ and its transpose.  Retain the explicit
flag-cover indices in its certificate.  On each of its eight open first
projective lines, project the used-variable factor to $k^7$ by the
integral matrix in the checker and apply the degree-three Koszul map.
For indices $2,3,9,12,13$ the generic rank is 343, above
$17\binom63=340$.  Exact determinant-polynomial factorization leaves
only $t=-1$ for indices $2,3,9,13$, and only $t=1$ for index 12;
infinity is excluded.  Full rational invariant-flag continuation at
these points excludes every branch.  The remaining first projective
families are indices 0 and 1 at flag depth two, and index 5 at depth three.

\paragraph{A specialization-safe continuation lemma.}
Let an invariant partial flag depend rationally on a parameter $t$.
All its common-root kernels and rank tests may be computed over
$K=\mathbb Q(t)$.  Record every nonconstant numerator and denominator
of every nonzero Gaussian pivot.  Also record the denominators of
rational-function coefficients introduced in the next steps and the
nonzero leading coefficients of any determinant-polynomial factorization.
There is then an explicitly known nonzero polynomial $D(t)$ such that,
at every finite $t_0$ with $D(t_0)\ne0$, all the row reductions, kernels,
and nonzero rank minors specialize validly.

To justify this directly, each Gaussian operation is addition,
multiplication, or division by a recorded pivot.  Outside its zero and
pole set it is an invertible elementary operation and preserves its
echelon pivot pattern.  The formulas for a kernel basis use the same
operations.  Therefore ranks and dimensions used to enumerate successor
lines do not change on the specified open set.  Coefficients of every
surviving rational-function vector are defined there.  No assertion is
made about the removed fibers; they are recomputed separately below.

If a common kernel over $K$ is two-dimensional, the next choice is
$\mathbb P^1_K$.  Introduce a second parameter $u$ and apply BR-005's
exact pencil-minor procedure over $K[u]$.  In every continuation used
here, an excluding map either removes the whole pencil or leaves
finitely many sections $u=u_i(t)$, together with a possible infinity
section.  All exceptional factors that need continuation split into
linear factors over $K$; an unsplit higher-degree factor would be
reported open.  The roots $u_i(t)$ and their pole sets are retained.
In particular, a two-parameter choice is never replaced by finitely
many sample points.  Factorization is an identity in $K[u]$ and the
recorded leading-coefficient and denominator conditions make its
specialization valid outside $D(t)=0$.

All resulting sections are inserted as exact invariant flag successors.
An inadmissible section is excluded by a genuine Koszul rank inequality;
all others are continued.  At a zero of $D$, restart the entire original
projective family at that parameter, instead of specializing a possibly
singular generic echelon basis.  The point at infinity is likewise
restarted separately.  This proves the completeness of the procedure.

\paragraph{Exact arithmetic of the projected maps.}
The projected map can lose its torus grading.  Its fixed rational columns
are first put in exact row-reduced form by SymPy's rational
\texttt{DomainMatrix}.  If its pivot coordinates are $P$ and the
remaining coordinates are $Q$, the quotient map is explicitly
\[
 v\longmapsto v_Q-v_P R_{P,Q},
\]
where $R$ is that row-reduced matrix.  All subsequent
$\mathbb Q(t)$ arithmetic occurs in this quotient.  The variable
columns and the next pencil are therefore computed exactly, not via
numerical elimination or an assumed generic rank.  SymPy's optional
integer/rational backend only accelerates these exact operations.

\paragraph{The finite family tree.}
The complete calculation has six one-parameter family nodes.  Three
are the first families 0, 1, and 5.  Two further families appear at
depths three and four in exceptional rational fibers of family 0,
and one appears at depth three in an exceptional rational fiber of
family 1.  Denote these children by $0a,0b,1a$, respectively.
They are specified by full preceding flags and kernel bases in the
consolidated certificate.  The following table gives the exact
exceptional polynomials; infinity is added in every row.
\[
\begin{array}{c|r|l|l}
\text{family}&\text{initial depth}&\text{linear exceptional factors}
    &\text{irreducible quadratic factor}\\\hline
0&2&t,\ t+1,\ t+\frac12,\ t+\frac13,\ t-1,\ t-2&t^2+4t-1\\
1&2&t,\ t+1,\ t+\frac12,\ t+\frac13,\ t+3,\ t-1,\ t-2&t^2+\frac13\\
5&3&t,\ t+1,\ t+\frac12,\ t-1,\ t-2&\text{none}\\
0a&3&t,\ t+1,\ t+\frac12,\ t+\frac13,\ t+3,\ t-1&t^2+\frac13\\
0b&4&t,\ t+1,\ t-1&\text{none}\\
1a&3&t,\ t+1,\ t+\frac12,\ t+\frac13,\ t-1,\ t-2&t^2+4t-1
\end{array}
\]
The exact pivot choices can give a larger exceptional set in another
valid run; the checker computes and covers the set it actually produces.
It does not trust this table as an exhaustive input list.

Over $\mathbb Q(t)$, each of these six nodes has no unhandled new
parameter branch and no full surviving extension.  All linear-factor
roots and infinity are recomputed over $\mathbb Q$.  Their finite
flag trees also end without a full survivor.  The only additional
projective branches they expose are precisely the three children
already listed.  Every child has strictly greater flag depth than
its parent, a condition the checker asserts.

For each quadratic factor $f$, the entire original family is restarted
over the exact number field $\mathbb Q[t]/(f)$, not at a floating-point
root.  Both conjugate geometric roots have the same rank and kernel
behavior under the corresponding field isomorphisms.  Each of the four
such number-field computations ends without a surviving extension or
an unhandled projective kernel.  This covers the nonreal roots of
$t^2+\frac13$ as well as the real roots of $t^2+4t-1$.

Consequently none of the three first families contains a viable
seventeen-plane.  Together with the five earlier projected-pencil
exclusions, the full incidence for $S_\diamond$ is empty.
Transposition gives the same result for its partner.  BR-005 therefore
extends from 57 excluded Borel ideals to all 59.  The Borel-fixed
Grassmannian reduction for determinant quadratic spaces gives
$\mu_{10}\ge18$.  Finally, the multidimensional border-substitution
inequality with a 36-dimensional quadratic image yields
\[
 36-10+18=44.
\]
The inclusion of the Veronese variety in the corresponding Segre
variety gives the asserted symmetric-border-rank consequence.

\paragraph{Consolidated checker and proof status.}
After the BR-004 and BR-005 commands, run
\texttt{python scripts/check\_det10\_last\_case.py}.
It reconstructs the eight projected pencils, the generic family
computations, every exceptional rational and number-field fiber,
and the family tree exposed by those fibers.  It records a family
as open whenever a kernel, algebraic factor, or full extension has
not been excluded.  It asserts the strict depth increase of each
new special-fiber family.  A final \texttt{PASS} is produced only
when the queue is empty and no such unresolved branch remains.
The main output is \texttt{certificates/det10\_last\_case.json}.
The three checker commands form a complete independent reconstruction
from signed minor formulas; they do not import inherited generators
or reported ranks.

This is again a proper Grassmannian/flag incidence with determinantal
conditions.  Restricting it to the parameter curves and their open
subsets gives the generic calculation; the removed closed fibers are
handled by exact base change.  This explicitly avoids inferring
emptiness of special fibers from emptiness of a generic fiber.
The EGA/FGA/SGA dictionary above applies.  No purity theorem,
$\ell$-adic weight assertion, or motivic principle is used.

\subsection*{BR-007: an exact nineteen-term determinant ten-plane curve}
\textbf{Status: proved; root complete proof and primary-source review,
and independent exact BigInt coefficient verification, completed.}

Over every algebraically closed field of characteristic zero,
\[
 \mu_{10}(U^{\mathrm{det}}_{4,2})\le19.
\]
The witness is the coordinate ten-plane
$S_\diamond=\{0,1,2,3,6,7,9,18,19,21\}$ used above.
Combining this upper construction with BR-006 gives
$\mu_{10}\in\{18,19\}$.  Thus $j=10$ is the first index for which
$\mu_j-j\ge8$, using the previously settled values through $j=9$.
The exact tenth minimum remains unresolved.  A gap of nine at $j=10$
would require excluding border rank eighteen for every one of the
fifty-nine determinant Borel types, rather than only the last
rank-seventeen survivor.

\paragraph{A tangent insertion lemma.}
Suppose a rank-$r$ curve contains a summand
$a_i(t)\otimes b_i(t)\otimes c_i(t)$, with
$b_i(t),c_i(t)\longrightarrow x$.  Introduce an independent first-factor
vector $a_\ast$ and a new last-factor coordinate $v$.  Replace the summand
by the following two pure terms:
\[
 (a_i(t)-t^{-1}a_\ast)\otimes b_i(t)\otimes c_i(t)
 +t^{-1}a_\ast\otimes(b_i(t)+tv)\otimes(c_i(t)+tv).
\]
Their difference from the original summand is exactly
\[
 a_\ast\otimes\bigl(b_i(t)\otimes v+v\otimes c_i(t)
                      +t v\otimes v\bigr).
\]
Thus one additional pure term adds
$a_\ast\otimes(x\otimes v+v\otimes x)$ to the limit.
Reversing the signs of the $a_\ast$ coefficients gives the negative
insertion.  This is a Laurent-polynomial identity, not an infinitesimal
approximation used without a limiting curve.

\paragraph{The two limiting points.}
Use the exact seventeen-term Waring curve for the $3\times3$ determinant
displayed in Conner--Huang--Landsberg, \emph{Bad and good news for
Strassen's laser method}, \S8, printed page 23 of the 6 May 2021 version.
Among its displayed matrix vectors are
\[
 m_1(t)=-t^{-5}e_{00}-e_{11}+t^5e_{22},\qquad
 m_{10}(t)=-t^{-4}e_{01}-t^{-1}e_{10}-t^5e_{22}.
\]
Transpose all matrix coordinates, which preserves the determinant tensor.
In the first of these terms, multiply the second and third factors by
$-t^5$ and the first factor by $t^{-10}$.  In the second, use
$-t^4,-t^4,t^{-8}$.  These operations leave each pure summand unchanged.
The resulting second and third factor vectors are
\[
 X(t)=e_{00}+t^5e_{11}-t^{10}e_{22}\longrightarrow e_{00},
 \qquad
 Y(t)=e_{10}+t^3e_{01}+t^9e_{22}\longrightarrow e_{10}.
\]
Apply the positive insertion at $X$ with $v=e_{13}$ and the negative
insertion at $Y$ with $u=e_{03}$, using the same $a_\ast$.
The resulting nineteen-term curve has limit
\[
 T_{\det_3}+a_\ast\otimes
 (e_{00}\otimes e_{13}+e_{13}\otimes e_{00}
       -e_{10}\otimes e_{03}-e_{03}\otimes e_{10}).
\]
The first nine contraction slices are the nine minors of the upper-left
$3\times3$ block.  The new slice is the minor with row pair $01$ and
column pair $03$.  These are exactly the ten minors in $S_\diamond$.

\paragraph{Integral cofactor normalization.}
The first-factor basis in a determinant tensor is indexed by the missing
row and column of the cofactor.  Apply its sign and complement-pair
reindexing to identify it with the ordered integral minor basis.
For $S_\diamond$ in increasing order, the map from row-major
$3\times3$ coordinates is
\[
 (9,+1),(8,-1),(7,+1),(6,-1),(5,+1),(4,-1),
 (3,+1),(1,-1),(0,+1),
\]
where each pair is (target index, multiplier); $a_\ast$ has target index
2 and remains independent.  If a convention identifies the polynomial
with a symmetrized tensor carrying a common factorial denominator, include
the corresponding nonzero scalar on these first nine coordinates.
It is invertible in characteristic zero and does not alter the new slice.
The literal coefficient check below shows that the displayed CHL curve
already has integral ordered coefficients $\pm1$ before this reindexing.

\paragraph{Complete exact coefficient verification.}
The independent checker is
\texttt{python scripts/check\_det10\_upper19.py}; it uses only the standard
library and generates \texttt{certificates/det10\_upper19.json}.
It reconstructs all seventeen published base points, checks all 729
ordered coordinates of their tensor sum, builds the nineteen modified
pure terms, and compares every nonpositive Laurent coefficient with the
literal signed-minor tensor for $S_\diamond$.  It finds exactly 36
nonzero constant coordinates for the base determinant and exactly 40
for the ten-plane target; every negative coefficient is zero.
All nineteen factor vectors, exponents and coefficients are saved.

For exact arithmetic put $\zeta=e^{2\pi i/12}$, $h=\sqrt5$,
$\phi=(1+h)/2$, and $\gamma=5^{-1/6}$.
The two cube roots in the last two published vectors are
$\gamma\phi$ and $\gamma(\phi-1)$, because
$\phi^3=2+h$ and $(\phi-1)^3=h-2$.
Their common scalar satisfies $\gamma^3=h/5$.
After doubling those two vectors and multiplying the whole tensor by 40,
the entire check takes place in the free rank-eight integer module
\[
 \mathbb Z[\zeta,h]/(\zeta^4-\zeta^2+1,\ h^2-5).
\]
No floating-point root or fitted coefficient is used.  The root agent
independently reran the full construction with BigInt arithmetic at
01:05:15 UTC on 22 November 2026; all coefficients and the cofactor map
matched.  This is an upper degeneration, independent of the unfinished
review of the lower-bound certificate.

\paragraph{Scope of an unsuccessful eighteen-term search.}
A separate finite-field experiment considered all 90 pairs of published
points with distinct limiting supports, all insertion-order pairs from
$\{1,\ldots,20\}^2$, and every omitted-point eighteen-plane limit of each
nineteen-point curve.  This is 36,000 curves and 684,000 omission tests
over $\mathbb F_{1000081}$.  None contained all nine old minor directions
and a combination with both required new tangent components nonzero.
The ratio of the two components was allowed to vary.  This is a negative
result for that finite modular search, not a proof that no eighteen-term
curve exists in characteristic zero.
